\tocsection{\textbf{11/13/26 Lecture 32:} Homology of Cones and Spheres}

Since $H_q(K)$ is a function of $C_q(K)$, $C_{q + 1}(K)$, and $C_{q - 1}(K)$, we see $H_q(K)$ depends only on the \emph{$(q + 1)$-skeleton}, $K(q + 1)$, of $K$ (which are the simplices of dimension at most $q + 1$ of $K$).\\

If $L$ is a simplicial complex in $\R^n$, $CL$, the \emph{cone} on $L$, is a new simplicial complex in $\R^{n + 1}$ defined such that the simplices of $CL$ are
\[ \{\tt{simplices of } L\} \cup \{v = (0, \ldots, 0, 1) \in \R^{n + 1}\} \cup \{(v, v_0, \ldots, v_k) \col (v_0, \ldots, v_k) \in L\} \]
Where, when we say ``simplicies of $L$'', we are implicitly embedding $L$ into $\R^{n} \times \{0\}$. \\


We haven't proved this yet, but it turns out that the homology groups are a homotopy invariant. Thus, the reason we should believe the following theorem is true is because and $\abs{CL}$ is homotopy equivalent to a point.

\begin{theorem}[\textbf{Theorem 32.1}: Homology of a Cone]
    \[ H_q(CL) = \left\{\begin{matrix} \Z & \tt{if } q = 0 \\ 0 & \tt{if } q > 0 \end{matrix}\right. \]
\end{theorem}
\begin{proof}{Proof:}
    Let $K = CL$. Since $CL$ has one component, we see it holds for $q = 0$.\\

    Suppose $q > 0$. Define $d : C_q(K) \to C_{q + 1}(K)$ by
    \[ d\sigma = \left\{\begin{matrix} (vv_0 \cdots v_q) & \tt{if } \sigma = (v_0 \cdots v_q) \tt{ is a simplex in } L \\ 0 & \tt{otherwise} \end{matrix}\right. \]
    We want to show $\partial(d\sigma) = \sigma - d(\partial\sigma)$ for any oriented $q$-simplex.\\

    If $\sigma \in L$ with $\sigma = (v_0, \ldots, v_q)$,
    \begin{align*}
        \partial(d\sigma) &= \partial(vv_0 \cdots v_q)\\
        &= (v_0, \ldots, v_q) + \sum (-1)^{i + 1}(vv_0, \ldots, \widehat{v}_i, \ldots, v_q)\\
        &= \sigma - \sum (-1)^{i}d(v_0, \ldots, \widehat{v}_i, \ldots, v_q)\\
        &= \sigma - d(\partial\sigma)
    \end{align*}
    Now, if $\sigma \notin L$, write $\sigma = (vv_0, \ldots, v_{q - 1})$ for $(v_0, \ldots, v_{q - 1}) \in L$. Then,
    \[ d(\sigma) = 0 \implies \partial(d\sigma) = 0 \]
    and
    \begin{align*}
        d(\partial\sigma) &= d\left((v_0, \ldots, v_{q - 1}) + \sum (-1)^{i + 1}(v, v_0, \ldots, \widehat{v}_i, \ldots, v_{q - 1})\right)\\
        &= (vv_0 \cdots v_{q - 1}) + 0\\
        &= \sigma
    \end{align*}
    So $\sigma - d(\partial\sigma) = 0 = \partial(d\sigma)$.\\

    Extending linearly, we see that $\partial(dc) = c - d(\partial c)$ for any $c \in C_q(K)$. Now, if $z \in Z_q(K)$, so
    \[ \partial(dz) = z - d(\partial z) = z \]
    Thus, $z \in B_q$, which implies $Z_q = B_q$, and so $H_q(CL) = 0$.

    $\qed$
\end{proof}

Let $\Delta^n$ be the standard $n$-simplex in $\R^{n + 1}$, a simplicial complex, where
\[ \abs{\Delta^n} = \left\{x \in \R^{n + 1} \col x_i \geq 0,\ \sum x_i = 1\right\} \]
Let $K = \Delta^n$. If we let $\Sigma^n$ be the simplices on the boundary of $\Delta^{n + 1}$, then $\abs{\Sigma^n} = S^n$ up to homeomorphism.\\

For $n = 2$:
\fig{0.75}{Lecture32Fig1}{}
$\abs{\Sigma^2}$ is the union of the four triangles.

\begin{theorem}[\textbf{Theorem 32.2}: Homology of $S^n$]
    \[ H_q(\Sigma^n) = \left\{\begin{matrix} \Z & \tt{if } q = 0 \tt{ or } n \\ 0 & \tt{otherwise} \end{matrix}\right. \]
\end{theorem}
\begin{proof}{Proof:}
    The $n$-skeletons of $\Sigma^n$ and $\Delta^{n + 1}$ are the same. Thus, $H_q(\Sigma^n) = H_q(\Delta^{n + 1})$ if $q \leq n - 1$. Since $\Delta^{n + 1}$ is $C\Delta^n$, we see that
    \[ H_q(\Sigma^n) = H_q(\Delta^{n + 1}) = \left\{\begin{matrix} \Z & q = 0 \\ 0 & q > 0 \end{matrix}\right. \]
    for $q = 0, \ldots, n - 1$.\\

    Now, what about when $q = n$? We see
    \[ B_n(\Sigma^n) = \operatorname{Im}(\partial : C_{n + 1}(\Sigma^n) \to C_n(\Sigma^n)) = 0 \quad (\tt{as } C_{n + 1}(\Sigma^n) = 0) \]
    Thus,
    \[ H_n(\Sigma^n) = Z_n(\Sigma^n) \]
    Since $\Sigma^n$ and $\Delta^{n + 1}$ have the same $n$-skeleton, $Z_n(\Sigma^n) = Z_n(\Delta^{n + 1})$. But we know $H_n(\Delta^{n + 1}) = 0$, so $Z_n(\Delta^{n + 1}) = B_n(\Delta^{n + 1})$. Now
    \[ B_n(\Delta^{n + 1}) = \operatorname{Im}(\partial : C_{n + 1}(\Delta^{n + 1}) \to C_n(\Delta^{n + 1})) \]
    but as $C_{n + 1}(\Delta^{n + 1})$ is generated by the single $(n + 1)$-simplex $\sigma$ in $\Delta^{n + 1}$, $C_{n + 1}(\Delta^{n + 1}) = \Z\sigma$, so
    \[ B_n(\Delta^{n + 1}) = \Z(\partial\sigma) \neq 0 \]
    Thus,
    \[ H_{n}(\Sigma^{n}) = Z_n(\Sigma^n) = Z_n(\Delta^{n + 1}) = B_n(\Delta^{n + 1}) \cong \Z \]
    For higher degrees, we see $C_q(\Sigma^n) = 0$, so
    \[ H_q(\Sigma^n) = \left\{\begin{matrix} \Z & \tt{if } q = 0 \tt{ or } n \\ 0 & \tt{else} \end{matrix}\right. \]

    $\qed$
\end{proof}

Recall $H_0(K) = \Z^r$ where $r$ is the number of components. What is $H_1(K)$?\\

The definition of $H_1(K)$ looks sort of like $\pi_1(\abs{K}, v)$, so are they the same? The answer is no; $\pi_1$ doesn't have to be abelian when $H_1(K)$ does.\\

However, that is mostly all we need to fix. We will show $H_1$ is the ``abelianization'' of $\pi_1$. Before this, some group theory.

\begin{definition}[\textbf{Definition 32.3:} Commutator Subgroup and Abelianization]
    If $G$ is a group, the \emph{commutator subgroup} $[G, G]$ consists of products of the form
    \[ (g_1g_1'g_1^{-1}(g_1')^{-1}) \cdots (g_kg_k'g_k^{-1}(g_k')^{-1}) \]
    That is, it is generated by products of the form $aba^{-1}b^{-1}$.\\

    This turns out to be a normal subgroup of $G$, and $G/[G, G]$ is the \emph{abelianization} $G^{ab}$ of $G$.
\end{definition}

\begin{example}{\textbf{Example:} Free Groups}
    $G = F_r = \{\tt{words in } r \tt{ letters and their inverses}\}$. Then, $G^{ab} = \Z^r$.
\end{example}
